%! Unit Opener: Functions as Models — Unit 2, Lesson 2.0 — Student Packet Cover Sheet
\documentclass[10pt]{article}
\usepackage{atda-article}
\usepackage{atda-boxes}
\usepackage{ltablex}
\keepXColumns

\begin{document}

% Full-bleed royal banner
\begin{tikzpicture}[remember picture, overlay]
  \fill[royal] (current page.north west)
    rectangle ([yshift=-0.9in]current page.north east);
\end{tikzpicture}
\vspace{-0.8in}

\begin{center}
  {\color{white}\textbf{\LARGE Algebra, Trigonometry, and Data Analysis}}\\[4pt]
  {\color{mistmid}\large\textbf{Unit 2: Functions, Transformations, and Graph Analysis}}\\[6pt]
  {\color{white}\textbf{\large Lesson 2.0 \quad Unit Opener: Functions as Models}}
\end{center}

\vspace{0.22in}
\namedateperiod
\vspace{0.18in}

\begin{learningtargetbox}
\small
\begin{enumerate}[leftmargin=1.5em, itemsep=5pt]
  \item \ldots{} read a graph in context and name what it tells me: the \textbf{domain} and
        \textbf{range}, where the function is \textbf{increasing}, \textbf{decreasing}, or
        \textbf{constant}, and its highest and lowest points.
  \item \ldots{} recognize the three function families this unit studies --- \textbf{linear},
        \textbf{quadratic}, and \textbf{exponential} --- from a graph or a table, and name the
        \textbf{parent function} of each.
  \item \ldots{} use the algebra Unit 2 is built on: function notation, values a formula excludes,
        rate of change, and the intercepts of a line.
  \item \ldots{} name what each lesson of Unit 2 teaches and where that skill gets used later in
        this course.
\end{enumerate}
\end{learningtargetbox}

\vspace{0.14in}

\begin{tocbox}
\small
\renewcommand{\arraystretch}{1.5}
\begin{tabularx}{\linewidth}{@{} c l X r @{}}
\rowcolor{mist}
\textbf{\#} & \textbf{Component} & \textbf{Description} & \textbf{Score} \\
\midrule
1 & Warm-Up        & Unit 2 diagnostic --- five prerequisite skills & \blank{1.2cm} \\
2 & Guided Notes   & One trail graph, every question; the three families; the Unit 2 roadmap
                                                                    & \blank{1.2cm} \\
3 & Group Activity & Reading a day off a graph: a rideshare fare and a phone battery
                                                                    & \blank{1.2cm} \\
4 & Exit Ticket    & Domain and range from a graph, and why context limits a domain
                                                                    & \blank{1.2cm} \\
5 & Homework       & Fluency practice, a viral-video graph, and a look ahead & \blank{1.2cm} \\
\midrule
  & \multicolumn{2}{r}{\textbf{Total}} & \blank{1.2cm} \\
\end{tabularx}
\end{tocbox}

\vspace{0.14in}

\begin{remindbox}
\small
This unit is about \emph{reading} functions: \texttt{AFDA.AF.1a--b} (parent functions and
transformations), \texttt{AFDA.AF.1d, f} (moving between an equation and a graph),
\texttt{AFDA.AF.1c, e, g} (choosing and comparing models), and \texttt{AFDA.AF.2a--h} (domain and
range, intervals, maximums and minimums, intercepts, end behavior, and asymptotes). In this
course those characteristics are read \textbf{from a graph} --- so bring your eyes, not just your
algebra. Today's lesson introduces no new standard; it maps the unit and finds out where you are
starting.
\end{remindbox}

\end{document}
